\documentclass[10pt, aspectratio=169]{beamer}
\usepackage{amsmath, amssymb, amsthm}
\usepackage{xcolor}

\usetheme{Madrid}
\usecolortheme{whale}

% Define custom theorem environments
\newtheorem{intuition}{Intuition}
\newtheorem{motivation}{Motivation}

\title[Phase 5: The Climax]{Phase 5: The Climax -- Squashing and Achieving FHE}
\subtitle{Fully Homomorphic Encryption over the Integers}
\author{Paper Presentation}
\date{\today}

\begin{document}

\frame{\titlepage}

\begin{frame}{Outline: Phase 5}
    \tableofcontents
\end{frame}

\section{The Depth Problem}

\begin{frame}{The Crisis of Bootstrapping}
    \begin{motivation}
    We established in Phase 1 that to achieve Fully Homomorphic Encryption, our scheme must be \textbf{bootstrappable}. It must be able to homomorphically evaluate its own decryption circuit before the noise exceeds the $p/2$ boundary.
    \end{motivation}
    \pause
    
    \vspace{0.3cm}
    Let us look at the mathematics of our decryption equation:
    $$m' = (c \pmod p) \pmod 2$$
    \pause
    
    To evaluate this homomorphically as a boolean circuit, we must compute it over the integers:
    $$c \pmod p = c - p \lfloor c/p \rceil$$
\end{frame}

\begin{frame}{The Wall of Integer Division}
    \textbf{The Roadblock:}
    To evaluate $c - p \lfloor c/p \rceil$, the circuit must compute $\lfloor c/p \rceil$. This requires \textbf{integer division}. 
    \pause
    
    \vspace{0.3cm}
    When translated into a boolean polynomial circuit, division has a massive multiplicative depth. 
    \pause
    
    \vspace{0.3cm}
    Recall our permitted polynomial degree bound $d$ from Phase 2:
    $$d \le \frac{\eta - 4 - \log||\hat{f}||_1}{\rho' + 2}$$
    \pause
    
    \textbf{The Reality:} The standard division circuit's degree far exceeds this permitted noise budget. The decryption circuit is simply too deep. We must \textbf{"squash"} it.
\end{frame}

\section{Squashing the Decryption Circuit}

\begin{frame}{Squashing: Shifting the Burden}
    To make the scheme bootstrappable, we must reduce the multiplicative depth of the decryption circuit.
    \pause
    
    \vspace{0.3cm}
    \begin{intuition}[The Squashing Strategy]
    Division by $p$ is mathematically expensive. Multiplication is relatively cheap. We will shift the heavy computational lifting from the $\text{Decrypt}$ algorithm (which is strictly constrained by the noise budget) to the $\text{Evaluate}$ algorithm (which can be done in the clear by the server \textit{before} the final homomorphic step).
    \end{intuition}
    \pause
    
    \vspace{0.3cm}
    \textbf{New Parameters:}
    \begin{itemize}
        \item $\kappa$: Precision of rational numbers.
        \item $\Theta$: Size of a massive public set of rational numbers.
        \item $\theta$: The Hamming weight of a sparse secret subset.
    \end{itemize}
\end{frame}

\begin{frame}{Modifying Key Generation: The Hint}
    During KeyGen, alongside $p$ and $pk^*$, we generate a "hint" about the value of $1/p$.
    \pause

    \vspace{0.2cm}
    \textbf{1. The New Secret Key ($\vec{s}$):}
    Choose a random $\Theta$-bit vector $\vec{s} = \langle s_1, \dots, s_\Theta \rangle$ with a very low Hamming weight $\theta$. 
    Let $S = \{i : s_i = 1\}$. This sparse vector $\vec{s}$ becomes the new secret key.
    \pause

    \vspace{0.2cm}
    \textbf{2. The Public Rational Numbers ($\vec{y}$):}
    Generate a set $\vec{y} = \{y_1, \dots, y_\Theta\}$ of rational numbers in $[0, 2)$ with $\kappa$ bits of precision, such that the subset $S$ sums to approximately $1/p$:
    $$\sum_{i \in S} y_i \approx \frac{1}{p} \pmod 2$$
    Specifically, $\left[ \sum_{i \in S} y_i \right]_2 = \frac{1}{p} - \Delta_p$, where $|\Delta_p| < 2^{-\kappa}$.
\end{frame}

\begin{frame}{Modifying Evaluate and Decrypt}
    When the server finishes homomorphically evaluating a circuit, it holds a noisy ciphertext $c^*$. \textbf{Before returning it}, the server does post-processing in the clear.
    \pause

    \vspace{0.2cm}
    \textbf{Server Post-Processing:}
    For $i \in \{1, \dots, \Theta\}$, the server computes:
    $$z_i = [c^* \cdot y_i]_2$$
    keeping only $n = \lceil \log \theta \rceil + 3$ bits of precision. The server outputs $c^*$ and $\vec{z} = \langle z_1, \dots, z_\Theta \rangle$.
    \pause

    \vspace{0.2cm}
    \textbf{Deriving the Squashed Decryption:}
    We need to compute $c^*/p \pmod 2$.
    $$\frac{c^*}{p} \approx c^* \sum_{i \in S} y_i \pmod 2 = \sum_{i=1}^\Theta s_i (c^* y_i) \pmod 2 \approx \sum_{i=1}^\Theta s_i z_i \pmod 2$$
    \pause
    
    \textbf{The New Decrypt Algorithm:}
    $$m' = \left[ c^* - \left\lfloor \sum_{i=1}^\Theta s_i z_i \right\rceil \right]_2$$
\end{frame}

\section{The Shallow Circuit \& Bootstrapping Achieved}

\begin{frame}{The Shallow Decryption Circuit}
    Look at our new decryption equation: $m' = \left[ c^* - \lfloor \sum_{i=1}^\Theta s_i z_i \rceil \right]_2$.
    \pause
    
    \vspace{0.2cm}
    \textbf{Why is this a mathematical triumph?}
    \begin{itemize}
        \item The server has already provided the rational numbers $z_i$.
        \item The secret key bits $s_i$ are strictly $0$ or $1$.
        \item Our only job is to sum up numbers. Because $\vec{s}$ is sparse, only $\theta$ of them are non-zero!
    \end{itemize}
    \pause

    \vspace{0.2cm}
    \begin{lemma}[4.1: Elementary Symmetric Polynomials]
    We use the classic "three-for-two" trick to reduce the sum. The $i$-th bit of the Hamming weight of a binary vector $\vec{\sigma}$ can be expressed as a polynomial of degree exactly $2^i$ using the $2^i$-th elementary symmetric polynomial $e_{2^i}(\vec{\sigma})$.
    \end{lemma}
\end{frame}

\begin{frame}{Bounding the Degree}
    Because we only need $\lceil \log \theta \rceil + 3$ bits of precision after the binary point, the maximum degree required to compute the binary sum and rounding is strictly bounded by:
    $$\text{Degree} \le 64\theta \log^2\theta$$
    \pause
    
    \vspace{0.3cm}
    \begin{intuition}
    By forcing the evaluator to pre-compute the partial products $z_i$ using the hint $y_i$, we squashed the massively deep integer division circuit down to a simple, sparse subset sum. 
    \end{intuition}
\end{frame}

\begin{frame}{The Grand Finale: Bootstrapping Achieved}
    By setting $\theta = \lambda$, the degree of our squashed decryption polynomial becomes $\mathcal{O}(\lambda \log^2 \lambda)$.
    \pause
    
    \vspace{0.3cm}
    \begin{theorem}[6.2: Bootstrapping Achieved]
    If we set our secret key length $\eta = \rho \cdot \Theta(\lambda \log^2 \lambda)$, the base SHE scheme's noise budget can comfortably handle polynomials of degree $\mathcal{O}(\lambda \log^2 \lambda)$. 
    \pause
    
    Because the degree of the squashed decryption circuit is strictly less than the maximum degree of the permitted polynomials, we have:
    $$\mathcal{D}_{\mathcal{E}} \subset \mathcal{C}_{\mathcal{E}}$$
    Therefore, the scheme is bootstrappable.
    \end{theorem}
    \pause
    
    \vspace{0.3cm}
    \begin{center}
        \textbf{We have achieved Fully Homomorphic Encryption using nothing but elementary modular arithmetic over the integers.}
    \end{center}
\end{frame}

\section{References}
\begin{frame}{References}
    \begin{thebibliography}{9}
        \bibitem{DGHV10}
        Marten van Dijk, Craig Gentry, Shai Halevi, and Vinod Vaikuntanathan.
        \textit{Fully Homomorphic Encryption over the Integers}. 
        Eurocrypt, 2010.
        
        \bibitem{Gentry09}
        Craig Gentry.
        \textit{A Fully Homomorphic Encryption Scheme}. 
        PhD Thesis, Stanford University, 2009.
        
        \bibitem{Boyar00}
        J. Boyar, R. Peralta, and D. Pochuev.
        \textit{On the multiplicative complexity of boolean functions over the basis ($\land$, $\oplus$, 1)}.
        Theor. Comput. Sci., 2000.
    \end{thebibliography}
\end{frame}

\end{document}